% !TeX root = elegantnote-cn.tex
\section{可逆电池的电动势及其应用}

\begin{enumerate}
    \item 等温等压条件下，系统发生变化时， Gibbs 自由能的减少等于对外所做的最大非膨胀功。
    \begin{equation*}
        (\Delta_{r}G)_{T,p} = W_{\mathrm{f,max}}
    \end{equation*}
\end{enumerate}

\subsection{可逆电池和可逆电极 20260316}

\begin{enumerate}
    \item 化学反应转变为产生电能的电池的条件：氧化还原反应，用装置使化学反应分别通过电极上的反应来完成。
    \item 必须有两个电极和相应电解质。
    \item 单液电池：电极插在同一种电解质；
    %pic
    双液电池：电极插在不同电解质。
    
    Daniell电池也是可逆电池（图9.1b），在电解 Daniell 电池时，阴极（锌极）液中虽然 \ce{H^+} 的热力学电势高于 \ce{Zn^{2+}}，但由于 \ce{H_2} 在锌表面的析氢超电势极大，使得 \ce{H^+} 实际上很难被还原。因此，\ce{Zn^{2+}} 能够避开 \ce{H^+} 的竞争，优先在阴极得到电子还原为 \ce{Zn} 金属。

    盐桥电池：要求阳离子和阴离子的迁移数尽可能相等。
\end{enumerate}

\subsubsection{可逆电池}

\begin{enumerate}
    \item 构成可逆电池的必要条件：
    \begin{enumerate}
        \item 电极上的化学反应可向正、反两个方向进行。
        \item 可逆电池存工作时，不论是充电还是放电，所通过的电流必须十分微小，电池是在接近平衡状态情况下工作的 。
    \end{enumerate}
\end{enumerate}

\subsubsection{可逆电极和电极反应}

\begin{enumerate}
    \item 第一类电极：金属与其阳离子组成的电极、氢电极、氧电极、卤素电极和汞齐电极。
    \item 第二类电极
    \item 第三类电极
    \item 设计电极（回放）
\end{enumerate}

\subsection{电动势的测定 20260318}

\subsubsection{对消法测电动势}

\begin{enumerate}
    \item Poffendorff 对消法：在外电路上加一个方向相反而电动势几乎相同的电池，以对抗原电池的电动势。
    \item 此时外电路几乎没有电流通过，相当于 $ R_{\text{o}} $ 无限大情况下进行测定。
\end{enumerate}

\subsubsection{标准电池}

\begin{enumerate}
    \item 标准电池电动势已知且稳定不变。
    \item Weston 标准电池：负极为镉汞齐，……
    \item 电池作用时的反应：
    
    负极

    正极

    净反应
    
    \item 标准电动势与温度的关系：
    \item 问题：
\end{enumerate}

\subsection{可逆电池的书写方法及电动势的取号 20260316 20260318}

\subsubsection{可逆电池的书写方法}

\begin{enumerate}
    \item 左边为负极，起氧化作用，是阳极；右边为正极，起还原作用，是阴极。
    \item | 表示相界面，有电势差存在。┆ 表示半透膜。%虚线
    \item || 或 ┆┆ 表示盐桥，使溶液见的接界电势通过盐桥降低到忽略不计。
\end{enumerate}

\subsubsection{从化学反应设计电池}

\begin{enumerate}
    \item 常规氧化还原反应
    \item 沉淀反应
    \item 中和反应 %equation
    
    \ce{Pt} | \ce{H2(g)} | \ce{H+} or \ce{OH-}  

    \ce{Pt} | \ce{O2(g)} | \ce{H+} or \ce{OH-} \\
    \ce{H+} 阴极氧气的电子变为水，阳极水失电子变为氧气。\\
    \ce{OH-} 阴极氧气的电子变为氢氧根，阳极氢氧根失电子变为水。

    \item Oxygen Reduction Reaction (ORR) 氧气还原，例如甲醇燃料电池中氧气在阴极被还原。
\end{enumerate}

\subsubsection{可逆电池电动势的取号}

\begin{enumerate}
    \item
    \begin{equation*}
        \Delta_{\text{r}}G_{\text{m}} = -zEF
    \end{equation*}
    \item 自发电池 $ \Delta_{\text{r}}G_{\text{m}} < 0 $ 非自发电池
    \item $ E $ 为强度性质
\end{enumerate}

\subsection{可逆电池的热力学 20260318}

\subsubsection{Nernst 方程}

\begin{enumerate}
    \item 
\end{enumerate}

\subsubsection{由标准电动势 $ E^{\circ} $ 求电池反应的平衡常数}

\subsubsection{由电动势 $ E $ 及其温度系数求反应的 $ \Delta_{\text{r}}H_{\text{m}} $ 和 $ \Delta_{\text{r}}S_{\text{m}} $}

\begin{enumerate}
    \item 根据热力学基本公式
    \begin{equation*}
        \begin{aligned}
            \mathrm{d}G = -S\mathrm{d}T + V\mathrm{d}p \\
            (\frac{\partial G}{\partial T})_{p} = -S
        \end{aligned}
    \end{equation*}
\end{enumerate}

\subsection{电动势产生的机理 20260318}

一个电池总的电动势可由下列三种电势差构成。

\subsubsection{电极与电解质溶液界面间电势差的形成}

\begin{enumerate}
    \item 在金属与溶液的界面上，由于正、负离子静电吸引和热运动两种效应的结果，溶液中的反离子只有一部分紧密地排在固体表面附近，相距约一、二个离子厚度称为紧密层；另一部分离子按一定的浓度梯度扩散到本体溶液中，称为扩散层。
    \item 紧密层和扩散层构成了双电层。
    \item 金属表面与溶液本体之间的电势差即为界面电势差。
\end{enumerate}

\subsubsection{接触电势}

\begin{enumerate}
    \item 逸出功的大小既与金属材料有关，又与金属的表面状态有关。
    \item 不同金属相互接触时，由于电子的逸出功不同，相互渗入的电子不同，在界面上电子分布不均匀，由此产生的电势差称为接触电势。
\end{enumerate}

\subsubsection{液体接界电势}

\begin{enumerate}
    \item 在两个含不同溶质的溶液的界面上，或溶质相同而浓度不同的界面上，由于离子迁移的速率不同而产生的电势差。
    \item 液接电势很小，一般在 \qty{0.03}{\volt}以下。
    \item 离子扩散是不可逆的，所以有液接电势存在的电池也是不可逆的，且液接电势的值很不稳定。
    \item 用盐桥可以使液接电势降到可以忽略不计。
\end{enumerate}

\subsubsection{}